

\section{Details of Proofs}
The appendix aims to introduce the complete proof of the previous lemmas and theorems. 
\subsection{The convergence of PI}

\begin{lemma}
Under update algorithm RWPI, the average reward should be monotonically increased.
\begin{proof}
From Algorithm 2, we could deduce the update rule of the average reward:
\begin{equation}
\begin{aligned}
J_{\pi_{t}}(\theta) - J_{\pi_{t-1}}(\theta) &= (W_{t}-1)J_{t-1}(\theta) + (1-W_{t})J_{\pi^{*}}(s,\theta)\\
&= (1-W_{t})(J_{\pi^{*}}(s,\theta)-J_{\pi_{t-1}}(\theta)).
\end{aligned}
\end{equation}
When $J_{\pi^{*}}(s,\theta) \geq J_{\pi_{t-1}}(\theta)$, we could deduce that $\log\frac{J_{\pi_{t}}(\theta)}{J_{\pi^{*}}(s,\theta)} \leq 1$. So the posterior weight $W_{t}$ is less than 1. This result holds vice versa. The first term $1-W_{t} \leq 0$ when $J_{\pi^{*}}(s,\theta) \leq J_{\pi_{t-1}}(\theta)$. Therefore, we could prove that:
\begin{equation}
J_{\pi_{t}}(\theta) - J_{\pi_{t-1}}(\theta) = (1-W_{t})(J_{\pi^{*}}(s,\theta)-J_{\pi_{t-1}}(\theta)) \geq 0.
\end{equation}
The sequence $J_{\pi_{t}}(\theta)$ is monotonically increased with time step $t$.
\end{proof}
\end{lemma}



\begin{lemma}
\label{appendix2}
Suppose Assumption 1 and Assumption 2 hold for some stochastic MDP $M$. Let $u_{i}$ be the state value in iteration $i$. Define $N$ as the maximum iteration number of the algorithm. Then $\pi_{t_{k}}$ is $\epsilon$-optimal after $N$ iterations.
\begin{proof}
Define $D = \min_{s}\{\mu_{i+1}(\pi)-\mu_{\pi}\}$ and $U = \max_{s}\{\mu_{i+1}(\pi)-\mu_{i}(\pi)\}$. Then we could deduce:
\begin{equation}
\begin{aligned}
D + \mu_{N}(\pi) &\leq \mu_{N+1}\\
&\leq W_{i}\mu_{N} + (1-W_{i})\pi_{i}^{*}(s,\theta)\\
&\leq W_{i}\mu_{N} + (1-W_{i})(r_{N} + \theta_{} u_{N}).
\end{aligned}
\end{equation}
Since $0 < W_{i} \leq 1$, the upper equation could be turned to:
\begin{equation}
D \leq (1-W_{i})J_{\pi_{i}}(\theta).
\end{equation}
Based on the definition in Preliminaries, let $\pi^{\star}$ be the optimal policy under all states that satisfies $\pi^{\star} := \sum_{s \in S}\pi_{i}^{\star}(s,\theta)$. Then
\begin{equation}
D \leq (1-W_{i})J_{\pi_{i}}(\theta) \leq (1-W_{i})J_{\pi^{\star}}(\theta).
\end{equation}
In a similar way, we could also prove $U \geq (1-W_{i})J_{\pi^{\star}}(\theta)$. From the definition of the stopping criterion of the Policy Iteration algorithm, we could assume $U - D \leq (1-W_{i})\epsilon$. Therefore, we have
\begin{equation}
\begin{aligned}
U &\leq D + (1-W_{i}) \\
&\leq (1-W_{i})J_{\pi_{i}}(\theta)+(1-W_{i})\epsilon\\
&\leq (1-W_{i})(J_{\pi_{i}}(\theta)+\epsilon)\\
(1-W_{i})J_{\pi^{\star}} &\leq (1-W_{i})(J_{\pi_{i}}(\theta)+\epsilon)\\
J_{\pi^{\star}} &\leq J_{\pi_{i}}(\theta)+\epsilon.
\end{aligned}
\end{equation}
We could deduce that stationary policy $\pi$ is $\epsilon$-optimal after $N$ iterations.
\end{proof}
\end{lemma}


\subsection{Regret Bound Analysis}

\begin{lemma}
$$
\log \frac{J_{\pi_{*}}(\theta)}{J_{\pi_{t}}(\theta)} \leq \frac{\epsilon}{C}
$$
\begin{proof}
First, we could multiply $J_{\pi_{t}}(\theta)$ in order to construct the inequality. Let $J_{\pi_{t}}(\theta) = n$, $\epsilon = x$
\begin{equation}
\begin{aligned}
\lim _{n \rightarrow+\infty}\left(1+\frac{x}{n}\right)^{n} &= \lim _{n \rightarrow+\infty} e^{n \ln \left(1+\frac{x}{n}\right)} \\
&=e^{\lim _{n \rightarrow+\infty} \frac{\ln \left(1+\frac{x}{n}\right)}{\frac{1}{n}}}.
\end{aligned}
\end{equation}
Apply the L'Hopital's Rule:
\begin{equation}
\begin{aligned}
\lim _{n \rightarrow+\infty}\left(1+\frac{x}{n}\right)^{n}&= e^{\lim _{n \rightarrow+\infty} \frac{\left(\frac{-x}{n^{2}}\right) \frac{1}{1+\frac{x}{n}}}{-\frac{1}{n^{2}}}}\\
&= e^{\lim _{n \rightarrow+\infty} \frac{x}{1+\frac{x}{n}}}=e^{x}.
\end{aligned}
\end{equation}
Then, we could prove that $\left(1+\frac{x}{n}\right)^{n}$ is monotonically increased with $n$:
\begin{equation}
\begin{aligned}
(1+\frac{x}{n})^{2} &= 1 \cdot\underbrace{ \left(1+\frac{x}{n}\right) \cdot\left(1+\frac{x}{n}\right) \cdots \cdots\left(1+\frac{x}{n}\right)}_{n}\\
& \leq \left[\frac{1+(1+\frac{x}{n})+ \cdots +(1+\frac{x}{n})}{n+1}\right]^{n+1}\\
& = \left[\frac{1+n(1+\frac{x}{n})}{n+1}\right]^{n+1}\\
& = \left[1+\frac{x}{n(n+1)}\right]^{n+1}\\
& \leq \left[1+\frac{x}{n+1}\right]^{n+1}.
\end{aligned}
\end{equation}
The first inequality holds for the arithmetic mean equality. We could deduce that $(1+\frac{x}{n})^{n} \leq e^{x}$. Therefore, we have:
\begin{equation}
J_{\pi_{t}}(\theta) \log \frac{J_{\pi_{*}}(\theta)}{J_{\pi_{t}}(\theta)} \leq \epsilon.
\end{equation}
Based on the definition of $C$, we could deduce the upper bound of average reward. Then the lemma could be proved.
\end{proof}
\end{lemma}



\begin{lemma}
\label{lemma3replacement}
Under Assumption 3, for each stationary near-optimal policy $\pi$ and epoch counter $k \geq 1$. Let $\rho(x)$ satisfies $\rho(x) := O(\sqrt{\log (\log(x))})$. The following upper bound holds for negative log-density.
$$
-\log W_{t_{k}}(\pi) \leq \frac{\epsilon}{C}|S|^{2}(  \rho(k_{\pi}) \sqrt{k_{\pi}} + k_{\pi} \tilde{\tau}_{t_{k},k_{\pi}})
$$
\begin{proof}
When $W_{t_{k}} \leq 1$, we could have:
\begin{equation}
W_{t_{k}}(\theta):=\exp \sum_{\pi, s_{1}, s_{2}} \mathcal{H}\left(N_{\pi}(k), \pi\right) \log \frac{J_{\pi_{t}}(\theta)}{J_{\pi_{*}}(\theta)}.
\end{equation}
Based on the definition of the counter $\mathcal{H}$, we could deduce the value of the posterior weight in a single epoch:
\begin{equation}
\begin{aligned}
W_{t_{k}}(\theta)&= \exp \left(\sum_{t=1}^{\infty} \mathbbm{1}\left\{\pi_{e(t)}=\pi,\left(S_{t}, S_{t+1}\right)=\left(s_{1}, s_{2}\right), N(e(t)) \leq k\right\}\log \frac{J_{\pi_{t}}(\theta)}{J_{\pi_{*}}(\theta)}\right)\\
& = \exp \left(\sum_{\pi \in \Pi} \sum_{\left(s_{1}, s_{2}\right) \in \mathcal{S}^{2}} \sum_{t=1}^{T} \mathbbm{1}\left\{\pi_{e(i)}=\pi,\left(S_{t}, S_{t+1}\right)=\left(s_{1}, s_{2}\right)\right\}\log \frac{J_{\pi_{t}}(\theta)}{J_{\pi_{*}}(\theta)}\right)\\
& = \exp \left(N_{\pi}(t) \sum_{\left(s_{1}, s_{2}\right) \in \mathcal{S}^{2}} \sum_{t=0}^{t-1} \frac{\mathbbm{1}\left\{\pi_{e(t)}=\pi,\left(S_{t}, S_{t+1}\right)=\left(s_{1}, s_{2}\right)\right\}}{N_{\pi}(t)}\log \frac{J_{\pi_{t}}(\theta)}{J_{\pi_{*}}(\theta)}\right).
\end{aligned}
\end{equation}
Where $N_{\pi}(t):=\sum_{t=0}^{t-1}\sum_{\pi \in \Pi} \mathbbm{1}\left\{\pi_{e(t)}=\pi\right\}$ represents the total number of the time instants during the period of $t$ when policy $\pi$ was conducted.

When Assumption 3 holds, we could know that $N_{\pi}(t) = \tilde{\tau}_{\pi_{t_{k}},N_{\pi}(k)}$, where $N_{\pi}(k) :=\sum_{k=0}^{K}\sum_{\pi \in \Pi} \mathbbm{1}\left\{\pi_{e(k)}=\pi\right\} $ holds for the number of the epochs where policy $\pi$ was chosen(The notation of $\tau$ will be represented as $N_{\pi}(k) = k_{\pi}$, $\tilde{\tau}_{\pi_{t_{k}},N_{\pi}(k)} = \tilde{\tau}_{t_{k},k_{\pi}}$). Therefore, we could have:
\begin{equation}
\begin{aligned}
& -\log W_{t_{k}}(\pi)\\
& = -N_{\pi}(t) \sum_{\left(s_{1}, s_{2}\right) \in \mathcal{S}^{2}} \sum_{t=0}^{t-1} \frac{\mathbbm{1}\left\{\pi_{e(t)}=\pi,\left(S_{t}, S_{t+1}\right)=\left(s_{1}, s_{2}\right)\right\}}{N_{\pi}(t)}\log \frac{J_{\pi_{t}}(\theta)}{J_{\pi_{*}}(\theta)}\\
& = -\sum_{\left(s_{1}, s_{2}\right) \in \mathcal{S}^{2}} \tilde{\tau}_{t_{k},k_{\pi}}  \mathcal{H}_{\left(s_{1}, s_{2}\right)}\left(\tilde{\tau}_{t_{k},k_{\pi}}, \pi\right) \log \frac{J_{\pi_{t}}(\theta)}{J_{\pi_{*}}(\theta)}\\
& = \sum_{\left(s_{1}, s_{2}\right) \in \mathcal{S}^{2}} \left[\tilde{\tau}_{t_{k},k_{\pi}}  \mathcal{H}_{\left(s_{1}, s_{2}\right)}\left(\tilde{\tau}_{t_{k},k_{\pi}}, \pi\right) - k_{\pi}\tilde{\tau}_{t_{k},k_{\pi}}\theta_{\pi}(s_{1}|s_{2})\right]\log \frac{J_{\pi_{*}}(\theta)}{J_{\pi_{t}}(\theta)} + \sum_{\left(s_{1}, s_{2}\right) \in \mathcal{S}^{2}} k_{\pi}\tilde{\tau}_{t_{k},k_{\pi}}\theta(s_{1}|s_{2}) \log \frac{J_{\pi_{*}}(\theta)}{J_{\pi_{t}}(\theta)}\\
\end{aligned}
\end{equation}
The last equation is based on the logarithmic property $\log \frac{A}{B} = - \log \frac{B}{A}$. Based on the Assumption 3, define $\rho(x) := O(\sqrt{\log (\log(x))})$.
\begin{equation}
\begin{aligned}
- \log W_{t_{k}}(\pi)& \leq \sum_{\left(s_{1}, s_{2}\right) \in \mathcal{S}^{2}} \rho(k_{\pi}) \sqrt{k_{\pi}} \log \frac{J_{\pi_{*}}(\theta)}{J_{\pi_{t}}(\theta)}+ k_{\pi}\tilde{\tau}_{t_{k},k_{\pi}} \sum_{\left(s_{1}, s_{2}\right) \in \mathcal{S}^{2}}\theta(s_{1}|s_{2}) \log \frac{J_{\pi_{*}}(\theta)}{J_{\pi_{t}}(\theta)}\\
& \leq \frac{\epsilon}{C}|S|^{2}(  \rho(k_{\pi}) \sqrt{k_{\pi}} + k_{\pi} \tilde{\tau}_{t_{k},k_{\pi}}).
\end{aligned}
\end{equation}
\end{proof}
\end{lemma}

\begin{lemma}
\label{lemma3.3replacement}
The difference between the local optimal average reward and the instantaneous average reward could be bounded by:
$$
|J^{*} - J_{\pi_{t}}| \leq \tilde{\mathcal{O}}(\frac{1}{\sqrt{T}}).
$$
\begin{proof}
We could know that the current policy probability distribution is updated based on the previous distribution and the current local optimal policy distribution:
\begin{equation}
\mu_{t}(\pi) = W_{t}\mu_{t-1}(\pi) + (1-W_{t})\pi_{t}^{*}(s,\theta_{t_{k}}).
\end{equation}
We could extend this result to reward function:
\begin{equation}
\label{equation50}
\begin{aligned}
J_{\pi_{t}}&=W_{t} J_{\pi_{t-1}}+\left(1-W_{t}\right) J^{*}\left(\theta_{t}\right)\\
J_{\pi_{t}}^{2} &=W_{t}^{2} J_{\pi{t-1}}^{2}+\left(1-W_{t}\right)^{2} J^{*^{2}}+2 W_{t}\left(1-W_{t}\right) J^{*} J_{\pi_{t-1}} \\
& \leq W_{t}^{2} J_{\pi t}^{2}+\left(1-W_{t}\right)^{2} J^{* 2}+2 W_{t}\left(1-W_{t}\right) J^{*} J_{\pi_{t}}.
\end{aligned}
\end{equation}
The inequality is based on the monotonicity of the algorithm. We could simplify Equation \ref{equation50}:
\begin{equation}
\begin{aligned}
\left(1-W_{t}^{2}\right) J_{\pi_{t}}^{2} &\leq \left(1-W_{t}\right)^{2} J^{*^{2}}+2 W_{t}\left(1-W_{t}\right) J^{*} J_{\pi_{t}} \\
\left(1+W_{t}\right) J_{\pi t}^{2} &\leq \left(1-W_{t}\right) J^{*^{2}}+2 W_{t} J^{*} J_{\pi_{t}} \\
J_{\pi_{t}}^{2}+W_{t} J_{\pi_{t}}^{2} &\leq J^{*^{2}}-W_{t} J^{*^{2}}+2 W_{t} J^{*} J_{\pi_{t}}\\
W_{t}\left(J_{\pi_{t}}^{2}+J^{*^{2}}\right) &\leq J^{*^{2}}-J_{\pi_{t}}^{2}+2 W_{t} J^{*} J_{\pi_{t}}\\
J_{\pi_{t}}^{2}+J^{*^{2}} &\leq \frac{1}{W_{t}}\left(J^{*^{2}}-J_{\pi_{t}}^{2}\right)+2 J^{*} J_{\pi_{t}}.
\end{aligned}
\end{equation}
Based on the definition of the regret of each time step, we could deduce the bound of the instantaneous regret:
\begin{equation}
\begin{aligned}
\left(J_{\pi t}-J^{*}\right)^{2} &=J_{\pi_{t}}^{2}+J^{*^{2}}-2 J_{\pi t} J^{*} \\
& \leq \frac{1}{W_{t}}\left(J^{*^{2}}-J_{\pi_{t}}^{2}\right)+2 J^{*} J_{\pi_{t}}-2 J_{\pi_{t}} J^{*} \\
&=\frac{1}{W_{t}}\left(J^{*^{2}}-J_{\pi_{t}}^{2}\right) \\
&=\frac{1}{W_{t}}\left(J^{*}-J_{\pi_{t}}\right)\left(J^{*}+J_{\pi_{t}}\right).
\end{aligned}
\end{equation}
Therefore, we could deduce that:
\begin{equation}
|J_{\pi_{t}}-J^{*}| \leq \frac{1}{W_{t}}|J_{\pi_{t}}+J^{*}|.
\end{equation}
From Lemma \ref{lemma3replacement}, we could know that $-\log W_{t_{k}}(\pi)$ is bounded by $B$, with $B = \frac{\epsilon}{C}|S|^{2}(\rho(k_{\pi}) \sqrt{k_{\pi}} + k_{\pi} \tilde{\tau}_{t_{k},k_{\pi}})$. Therefore, we could construct the following inequalities.
\begin{equation}
\begin{aligned}
W_{t_{k}}-1 &\geq \log W_{t_{k}} \geq -B\\
\frac{1}{W_{t_{k}}} &\leq \frac{1}{1-B}.
\end{aligned}
\end{equation}
Factor $B$ is proportional to parameter $k_{\pi}$ which could be bounded by the total number of episode of under total time $T$. Therefore, we could bound $\frac{1}{W_{t}}$ by $T$(Ignoring the constants):
\begin{equation}
\begin{aligned}
\frac{1}{W_{t}} &\leq \frac{1}{1-\frac{\epsilon}{C}|S|^{2}(\rho(k_{\pi}) \sqrt{k_{\pi}} + k_{\pi} \tilde{\tau}_{t_{k},k_{\pi}})}\\
&\leq \frac{1}{1-\sqrt{\sqrt{T}}-\sqrt{T}}.
\end{aligned}
\end{equation}
Based on the definition of $C$, the average reward function is bounded by $C$. So the difference between the local optimal average reward and the instantaneous average reward could be bounded by:
\begin{equation}
\begin{aligned}
|J^{*} - J_{\pi_{t}}| &= |J_{\pi_{t}}-J^{*}|\\
&\leq \frac{1}{W_{t}}|J^{*}+J_{\pi_{t}}|\\
&\leq \frac{2}{1-\frac{\epsilon}{C}|S|^{2}(\rho(k_{\pi}) \sqrt{k_{\pi}} + k_{\pi} \tilde{\tau}_{t_{k},k_{\pi}})}\\
&\leq \tilde{\mathcal{O}} (\frac{2C}{S^{2}\sqrt{T}}).
\end{aligned}
\end{equation}
\end{proof}
\end{lemma}

\begin{lemma}
\label{appendix4.7}

$$\frac{C}{S^{2}} < D\sqrt{SA}$$ when $|S| \geq 2$ or $|A| \geq 2$

\begin{proof}
We could know that $C$ is defined as the upper bound of the average reward. So we could deduct:
\begin{equation}
\begin{aligned}
&C \geq J_{\pi_{*}}\left(\theta_{*}\right)\\
&D \geq \max T_{s \rightarrow s^{\prime}}^{\pi t_{k}}\left(\theta_{t_{k}}\right)\\
&C \leq \max T_{s \rightarrow s^{\prime}}^{\pi_{*}}\left(\theta_{*}\right) \leq \max T_{s \rightarrow s^{\prime}}^{\pi t_{k}}\left(\theta_{t_{k}}\right) \leq D.
\end{aligned}
\end{equation}
Assuming Lemma \ref{appendix4.7} is established, we could get:
\begin{equation}
\begin{aligned}
&\frac{C}{S^{2}} \leq \frac{D}{S^{2}} \leq D\sqrt{SA}\\
&D \leq DS^{2}\sqrt{SA}\\
&S^{2}\sqrt{SA} \geq 1.
\end{aligned}
\end{equation}
Therefore the Lemma could be proved when the fMDP process has more than one state and one action. 
\end{proof}
\end{lemma}

\begin{theorem}
\label{theorem4appendix}
The first part of the regret in time step $T$ is bounded by:
$$
Reg_{T}^{1} \leq \tilde{\mathcal{O}}(\frac{\sqrt{T}}{S^{2}}).
$$
\begin{proof}
From the definition before, we could know that $Reg_{T}^{1}$ could be represented as:
\begin{equation}
Reg_{T}^{1} = TJ_{\pi_{{k}}}(\theta_{K}) - \sum_{t = 1}^{T}J_{\pi_{t}}(\theta_{K}).
\end{equation}
Since this theorem won't involve the transformation of the transition probability. So let $J_{\pi}(\theta) = J_{\pi}$. Based on the update rule of the posterior distribution $\mu_{t+1}(\pi)$ of policy $\pi$. We could divide the average reward into several parts:\\
At time step $t = T$, we could assume the instantaneous regret equals to zero:
\begin{equation}
Reg_{t_{T}}^{1} = J_{\pi_{ k}}- J_{\pi_{T}} = 0.
\end{equation}
At time step $t = T-1$, define the local optimal average reward as $J_{\pi}^{*}$. Note that this local optimal value is virtual. The instantaneous regret could be represented as:
\begin{equation}
\begin{aligned}
Reg_{t_{T-1}}^{1} &= J_{\pi_{ k}} - J_{\pi_{T-1}}\\
&= W_{t-1}J_{\pi_{T-1}} + (1-W_{t-1})J_{\pi}^{*} - J_{\pi_{T-1}}\\
&= (W_{t-1}-1)J_{\pi_{T-1}} + (1-W_{t-1})J_{\pi}^{*}\\
&= (1-W_{t-1})(J_{\pi}^{*} - J_{\pi_{T-1}}).
\end{aligned}
\end{equation}
In a similar fashion, at time step $t = T-2$, the instantaneous regret could be represented as:
\begin{equation}
\begin{aligned}
Reg_{t_{T-2}}^{1} &= J_{\pi_{ k}} - J_{\pi_{T-2}}\\
&= J_{\pi_{ k}} - J_{\pi_{T-1}} + J_{\pi_{T-1}} - J_{\pi_{T-2}}\\
&= (1-W_{t-1})(J_{\pi}^{*} - J_{\pi_{T-1}}) + (1-W_{t-2})(J_{\pi}^{*} - J_{\pi_{T-1}}).
\end{aligned}
\end{equation}
Based on Lemma \ref{lemma3.3replacement}, the difference between the local optimal value and the current average reward could be bounded by:
\begin{equation}
|J_{\pi}^{*} - J_{\pi_{t}}| \leq \tilde{\mathcal{O}}(\frac{1}{\sqrt{T}}).
\end{equation}

The sub-optimal models are sampled when their posterior probability is larger than $\frac{1}{T}$. This ensures the time complexity of the Thompson sampling process is no more than $O(1)$. So we could deduce the total regret in time step $T$.
\begin{equation}
\begin{aligned}
Reg_{T}^{1} &= \frac{1}{T} (Reg_{t_{T-1}}^{1} + Reg_{t_{T-2}}^{1} + \cdots + Reg_{t_{1}}^{1})\\
&\leq \tilde{\mathcal{O}}(\frac{2C}{S^{2}\sqrt{T}})(\frac{T-1}{T} + \frac{T-2}{T} + \cdots \frac{1}{T})\\
&\leq \tilde{\mathcal{O}}(\frac{2C \sqrt{T}}{S^{2}}).
\end{aligned}
\end{equation}
\end{proof}
\end{theorem}

In order to deduce the second regret bound generated by the transition probability, we should analyze our algorithm's performance over $T$ time step. We define the number of macro episodes $M = \mathbbm{1}\{t_{k} \leq T\}$. An episode is defined as the set of the time steps under stopping criterions. Therefore, we could deduce the bound of the number of episode.
\begin{lemma}
\label{lemma1.1}
Under the stopping criterion, the number of episodes $M$ could be bounded by:
$$
M \leq S A \log (T).
$$\cite{NIPS2017_36e729ec}
\begin{proof}
The stopping criterion is triggered whenever the visits number of the initial state-action pair is doubled. So $M$ could be represented as:
\begin{equation}
M_{(s, a)}=\left\{k \leq K_{T}: N_{t_{k}}(s, a)>2 N_{t_{k-1}}(s, a)\right\}.
\end{equation}
Since the number of the visit to state-action pair $(s,a)$ is doubled at the beginning of every epoch $k$. The size of $\mathcal{M}_{(s, a)}$ should be no larger than $O(\log (T))$. Assume $\left|\mathcal{M}_{(s, a)}\right| \geq \log \left(N_{T+1}(s, a)\right)+1$. We could have:
\begin{equation}
\begin{aligned}
N_{t_{K_{T}}}(s, a) &= \prod_{k \leq K_{T}, N_{t_{k-1}}(s, a) \geq 1} \frac{N_{t_{k}}(s, a)}{N_{t_{k-1}}(s, a)}\\
&> \prod_{k \in \mathcal{M}_{(s, a)}, N_{t_{k-1}}(s, a) \geq 1} 2 \\
&\geq N_{T+1}(s, a).
\end{aligned}
\end{equation}
This contradicts the fact that $N_{t_{K_{T}}}(s, a) \leq N_{T+1}(s, a)$. This leads to $\left|\mathcal{M}_{(s, a)}\right| \leq \log \left(N_{T+1}(s, a)\right)$. Therefore, we could obtain the bound of the number of the episodes:
\begin{equation}
\begin{aligned}
M & \leq 1+\sum_{(s, a)}\left|\mathcal{M}_{(s, a)}\right|\\
&\leq 1+\sum_{(s, a)} \log \left(N_{T+1}(s, a)\right) \\
& \leq 1+S A \log \left(\sum_{(s, a)} N_{T+1}(s, a) / S A\right)\\
&=1+S A \log (T / S A).
\end{aligned}
\end{equation}
Since the logarithmic function is concave, we could simplify the inequality to:
\begin{equation}
M \leq S A \log(T).
\end{equation}
\end{proof}
\end{lemma}

\begin{lemma}
\label{lemma1.2}
The total number of episodes of total time step $T$ could be bounded by:
$$
K_{T} \leq \sqrt{2 S A T \log (T)}.
$$\cite{NIPS2017_36e729ec}
\begin{proof}
Define macro episodes with start times $t_{n_{i}},i=1,2,\cdots$ where $t_{n_{1}}=t_{1}$,we could have
$$
t_{n_{i+1}}=\min \left\{t_{k}>t_{n_{i}}: N_{t_{k}}(s, a)>2 N_{t_{k-1}}(s, a) \right\}.
$$
Let $\tilde{T}_{i}=\sum_{k=n_{i}}^{n_{i+1}-1} T_{k}$ be the length of the ith episode. Therefore, within the $i$th macro episode, $T_{k} = T_{k-1}+1$ for all $k = n_{i},n_{i}+1,\cdots,n_{i+1}-2$.
\begin{equation}
\begin{aligned}
\tilde{T}_{i}&=\sum_{k=n_{i}}^{n_{i+1}-1} T_{k} \\ &=\sum_{j=1}^{n_{i+1}-n_{i}-1}\left(T_{n_{i}-1}+j\right)+T_{n_{i+1}-1} \\
& \geq \sum_{j=1}^{n_{i+1}-n_{i}-1}(j+1)+1=0.5\left(n_{i+1}-n_{i}\right)\left(n_{i+1}-n_{i}+1\right) .
\end{aligned}
\end{equation}
Consequently,$n_{i+1}-n_{i} \leq \sqrt{2 \tilde{T}_{i}}$, for all $i = 1,\cdots,M$. From this property, we could obtain:
\begin{equation}
\label{equation29}
K_{T}=n_{M+1}-1=\sum_{i=1}^{M}\left(n_{i+1}-n_{i}\right) \leq \sum_{i=1}^{M} \sqrt{2 \tilde{T}_{i}}.
\end{equation}
Based on Equation \ref{equation29} and $\sum_{i=1}^{M} \tilde{T}_{i}=T$, we could get:
\begin{equation}
K_{T} \leq \sum_{i=1}^{M} \sqrt{2 \tilde{T}_{i}} \leq \sqrt{M \sum_{i=1}^{M} 2 \tilde{T}_{i}}=\sqrt{2 M T}.
\end{equation}
Where the second inequality is based on Cauchy-Schwarz inequality. From Lemma \ref{lemma1.1}, we could know that the number of the macro episodes until time $T$ is bounded by $M \leq SA\log(T)$. Therefore, the lemma could be proved.
\end{proof}
\end{lemma}


\begin{theorem}
The regret caused by transition matrix sampling could be bounded by:
$$
Reg_{T}^{2} \leq \tilde{O}(D(2SAT)^{\frac{1}{4}}).
$$
\begin{proof}
In this theorem, we mainly focus on the difference between transition probability. From the previous definition of the Bellman iterator of the average reward, we could have:
\begin{equation}
\begin{aligned}
&J(\theta_{K})+b(\theta_{K}, \pi, s)=r(s, \pi)+\sum_{s^{\prime}} \theta_{K}\left(s^{\prime} \mid s, \pi\right) b\left(\theta_{K}, \pi, s^{\prime}\right) \\
&J\left(\theta_{t}\right)+b\left(\theta_{t}, \pi, s\right)=r(s, \pi)+\sum_{s^{\prime}} \theta_{t}\left(s^{\prime} \mid s, \pi\right) b\left(\theta_{t}, \pi, s^{\prime}\right).
\end{aligned}
\end{equation}
The difference between the average reward under near-optimal transition probability and instantaneous transition probability could be represented as:
\begin{equation}
\begin{aligned}
J(\theta_{K})-J\left(\theta_{t}\right) &=b\left(\theta_{t}, \pi\right)-b(\theta_{K}, \pi) \\
&+\sum_{s^{\prime},s} \left(\theta_{K}\left(s^{\prime} \mid s, \pi\right) b\left(\theta_{K}, \pi, s^{\prime}\right)-\theta_{t}\left(s^{\prime} \mid s, \pi\right) b\left(\theta_{t}, \pi, s^{\prime}\right)\right).
\end{aligned}
\end{equation}
We could bound the first term with the largest difference between each state:
\begin{equation}
\label{equation55}
\begin{aligned}
&0 \leq b(\theta_{K}, \pi, s)-b\left(\theta_{t}, \pi, s\right) \leq s p(b(\theta)) \leq D \\
&0 \geq b\left(\theta_{t}, \pi, s\right)-b(\theta_{K}, \pi, s) \geq-D.
\end{aligned}
\end{equation}
Based on Equation \ref{equation55}, we could bound the second term in a similar way:
\begin{equation}
\label{equation56}
\begin{aligned}
&\sum_{s^{\prime}} \theta_{K}\left(s^{\prime} \mid s, \pi\right) b\left(\theta_{K}, \pi, s^{\prime}\right)-\theta_{t}\left(s^{\prime} \mid s, \pi\right) b\left(\theta_{t}, \pi, s^{\prime}\right) \\
&\leq D \sum_{s^{\prime}}\left(\theta_{K}\left(s^{\prime} \mid s, \pi\right)-\theta_{t}\left(s^{\prime} \mid s, \pi\right)\right).
\end{aligned}
\end{equation}

Then, we define the total transition difference as:

\begin{equation}
\label{equation41}
\begin{aligned}
\theta_{*}(s^{\prime}|s,\pi) - \theta_{t}(s^{\prime}|s,\pi) &= \theta_{K}(s^{\prime}|s,\pi) - \theta_{t}(s^{\prime}|s,\pi)\\
\theta_{k}(s^{\prime}|s,\pi) - \theta_{t}(s^{\prime}|s,\pi) &= \theta_{K}(s^{\prime}|s,\pi) - \theta_{K-1}(s^{\prime}|s,\pi) + \theta_{K-1}(s^{\prime}|s,\pi) - \dots - \theta_{t}(s^{\prime}|s,\pi).
\end{aligned}
\end{equation}
From Equation \ref{equation41}, we could deduct the one-step transition difference to be:
\begin{equation}
\theta_{t}(s^{\prime}|s,\pi) - \theta_{t-1}(s^{\prime}|s,\pi) = \frac{\theta_{t-1}(s^{\prime}|s,\pi)+H\left(N_{\pi}(k), \pi\right)}{N_{\pi}(k)} - \theta_{t-1}(s^{\prime}|s,\pi).
\end{equation}
Based on Assumption 3, we could bound the one-step transition difference:
\begin{equation}
\begin{aligned}
\sum_{s_{1},s_{2}}(\theta_{t} - \theta_{t-1}) &= \sum_{s_{1},s_{2}}\left[\frac{\theta_{t-1}+H_{s_{1},s_{2}}}{N} - \theta_{t-1}\right]\\
&= \sum_{s_{1},s_{2}}\frac{\theta_{t-1}+H_{s_{1},s_{2}}}{N} - \sum_{s_{1},s_{2}}\theta_{t-1}\\
&= \sum_{s_{1},s_{2}}\frac{\theta_{t-1}+H_{s_{1},s_{2}}}{N} - T_{s_{1} \rightarrow s_{2}}^{\pi}(\theta_{t-1})\\
&= \sum_{s_{1},s_{2}}\frac{\theta_{t-1}+H_{s_{1},s_{2}}}{N} - \tilde{\tau}_{t-1}.
\end{aligned}
\end{equation}
where
\begin{equation}
\sum_{s_{1},s_{2}}\frac{\theta_{t-1}+H_{s_{1},s_{2}}}{N} \geq \frac{H_{s_{1},s_{2}(k,\pi)}}{N}.
\end{equation}
When we use stationary policy $\pi$ in epoch $k$, we could know that the count function of the policy $\pi$ should be less or equal to the total number of epochs. Therefore, we could deduct that:
\begin{equation}
 \sum_{s_{1},s_{2}}\frac{\theta_{t-1}+H_{s_{1},s_{2}}}{N} \geq \frac{H_{s_{1},s_{2}(k,\pi)}}{N} \geq \frac{H_{s_{1},s_{2}(k,\pi)}}{k}.
\end{equation}
Based on Assumption 3, we could deduct the bound for $\sum_{s_{1},s_{2}}(\theta_{t} - \theta_{t-1})$:

\begin{equation}
\label{equation46}
\sum_{s_{1},s_{2}}(\theta_{t} - \theta_{t-1}) \leq \sqrt{\frac{e_{1} \log \left(e_{2} \log k\right)}{k}}.
\end{equation}
Combining Equation \ref{equation46} with Equation \ref{equation41}. Since the update of the transition matrix only happens once in each epoch, we could deduct the difference between the periodic transition matrix and instantaneous transition matrix based on \ref{lemma1.2}:
\begin{equation}
\label{equation54}
\begin{aligned}
\theta_{K}\left(s^{\prime} \mid s, \pi\right)-\theta_{1}\left(s^{\prime} \mid s, \pi\right) &\leq K\sqrt{\frac{e_{1} \log \left(e_{2} \log k\right)}{k}} \\
&\leq \sqrt{k e_{1} \log \left(e_{2} \log k\right)}\\
&\leq \sqrt{\sqrt{2SAT\log{T}} e_{1} \log \left(e_{2} \log k\right)}.
\end{aligned}
\end{equation}

Therefore, we could combine Equation \ref{equation54} with Equation \ref{equation56}:

\begin{equation}
\begin{aligned}
\sum_{t=1}^{T}J(\theta_{K})-\sum_{t=1}^{T}J\left(\theta_{t}\right) &\leq D\sqrt{\sqrt{2SAT\log{T}} e_{1} \log \left(e_{2} \log k\right)}\\
&\leq \tilde{O}(D\sqrt{\sqrt{2SAT}}).
\end{aligned}
\end{equation}
\end{proof}
\end{theorem}


